\documentclass[11pt]{article}
\newcommand{\DocShortTitle}{Science-Team Rationale r0.5}
\newcommand{\DocPdfTitle}{SCI-CAL v0.1 TolTEC Timestream Calibration - Science-Team Rationale r0.5 Frozen}
\usepackage[letterpaper,margin=0.82in,headheight=14pt]{geometry}
\usepackage[T1]{fontenc}
\usepackage{lmodern}
\usepackage{amsmath,amssymb,mathtools}
\usepackage{booktabs,longtable,tabularx,array}
\usepackage{enumitem}
\usepackage{xcolor}
\usepackage{microtype}
\usepackage{fancyhdr}
\usepackage{hyperref}

\definecolor{calblue}{HTML}{244B6B}
\definecolor{calteal}{HTML}{287A78}
\definecolor{calgray}{HTML}{4D5963}
\definecolor{callight}{HTML}{EEF3F6}
\definecolor{calwarn}{HTML}{8A4B08}

\hypersetup{
  colorlinks=true,
  linkcolor=calblue,
  urlcolor=calteal,
  citecolor=calblue,
  pdftitle={\DocPdfTitle},
  pdfauthor={SCI-CAL independent scientific author},
  pdfsubject={SCI-CAL v0.1 frozen scientific contract}
}

\setlength{\emergencystretch}{3em}
\setlength{\parindent}{0pt}
\setlength{\parskip}{0.52em}
\setlength{\tabcolsep}{4pt}
\renewcommand{\arraystretch}{1.18}
\setlist{nosep,leftmargin=*}
\newcolumntype{L}[1]{>{\raggedright\arraybackslash}p{#1}}
\newcolumntype{Y}{>{\raggedright\arraybackslash}X}

\pagestyle{fancy}
\fancyhf{}
\fancyhead[L]{\small\color{calgray}SCI-CAL v0.1 FROZEN}
\fancyhead[R]{\small\color{calgray}\DocShortTitle}
\fancyfoot[C]{\small\thepage}
\renewcommand{\headrulewidth}{0.35pt}
\renewcommand{\footrulewidth}{0pt}

\newcommand{\code}[1]{\nolinkurl{#1}}
\newcommand{\Req}[1]{\texttt{SCI-CAL-REQ-#1}}
\newcommand{\Assump}[1]{\texttt{SCI-CAL-ASM-#1}}
\newcommand{\Edge}[1]{{\scriptsize\texttt{SCI-CAL-EDGE-#1}}}
\newcommand{\Term}[1]{\textbf{#1}}
\newcommand{\E}{\mathbb{E}}
\newcommand{\Cov}{\operatorname{Cov}}
\newcommand{\Var}{\operatorname{Var}}
\newcommand{\diag}{\operatorname{diag}}
\newcommand{\TOA}{\mathrm{TOA}}
\newcommand{\valid}{\mathcal{J}_{\mathrm{valid}}}
\newcommand{\mjybeam}{\mathrm{mJy\,beam^{-1}}}

\newcommand{\AuthorityNotice}{%
  \begin{center}
  \fcolorbox{calblue}{callight}{%
    \begin{minipage}{0.91\linewidth}
    \small\textbf{Frozen authority boundary.}
    This is an implementation-blind frozen scientific contract. It defines
    authoritative scientific meaning and falsifiable conformance conditions. It is
    not evidence of implementation conformity, atmosphere-model fidelity,
    observational repeatability, absolute flux accuracy, production readiness,
    or completed validation.
    \end{minipage}}
  \end{center}
}

\newcommand{\MakeContractTitle}[2]{%
  \hypersetup{pageanchor=false}%
  \begin{titlepage}
  \centering
  \vspace*{0.55in}
  {\Large\color{calteal}\textbf{SCI-CAL}}\\[0.15in]
  {\huge\color{calblue}\textbf{Detector Calibration, Atmospheric Extinction,\\[0.08in]
  and Signal Transfer}}\\[0.32in]
  {\LARGE\textbf{#1}}\\[0.22in]
  {\large Version 0.1 / \DocRevision{} - FROZEN}\\[0.12in]
  {\large \DocDate{}}\\[0.48in]
  \fcolorbox{calblue}{callight}{%
    \begin{minipage}{0.78\linewidth}
    \centering\small
    #2
    \end{minipage}}
  \vfill
  {\large Scientific owner: Grant Wilson}\\[0.08in]
  {\normalsize Independent, implementation-blind authorship}\par
  \end{titlepage}
  \hypersetup{pageanchor=true}%
}

\fancyhead[L]{\small\color{calgray}SCI-CAL v0.1 -- Science-Team Rationale r0.5 FROZEN}
\fancyhead[R]{}
\hypersetup{pdftitle={SCI-CAL v0.1 TolTEC Timestream Calibration - Science-Team Rationale r0.5 Frozen},pdfauthor={SCI-CAL scientific owner review revision},pdfsubject={SCI-CAL v0.1 frozen scientist-facing rationale r0.5}}

\newcommand{\DecisionBox}[2]{\begin{center}\fcolorbox{calteal}{callight}{\begin{minipage}{0.91\linewidth}\small\textbf{Scientific-owner decision: #1}\par #2\end{minipage}}\end{center}}
\newcommand{\StatusBox}[1]{\begin{center}\fcolorbox{calblue}{callight}{\begin{minipage}{0.91\linewidth}\small #1\end{minipage}}\end{center}}
\newcommand{\flowbox}[1]{\fcolorbox{calteal}{callight}{\parbox[c][0.48in][c]{0.20\linewidth}{\centering\small #1}}}
\newcommand{\flowboxsmall}[1]{\fcolorbox{calteal}{callight}{\parbox[c][0.72in][c]{0.145\linewidth}{\centering\scriptsize #1}}}

\begin{document}
\hypersetup{pageanchor=false}
\begin{titlepage}
\centering\vspace*{0.55in}
{\Large\color{calteal}\textbf{SCI-CAL}}\\[0.18in]
{\huge\color{calblue}\textbf{TolTEC Timestream Calibration}}\\[0.08in]
{\LARGE\color{calblue}\textbf{Scientific Basis, Assumptions, and Validation}}\\[0.38in]
{\large Scientific contract version 0.1}\\[0.06in]
{\large Rationale revision 0.5 -- FROZEN}\\[0.10in]
{\large 2026-08-23}\\[0.50in]
\fcolorbox{calblue}{callight}{\begin{minipage}{0.78\linewidth}\centering\small
A science-team explanation of the calibration chain, its physical meaning,
uncertainty boundaries, validity limits, and evidence still required. The separate
SCI-CAL engineering contract remains the normative conformance authority.
\end{minipage}}
\vfill {\large Scientific owner: Grant Wilson}
\end{titlepage}
\hypersetup{pageanchor=true}

\section*{Executive summary}
\addcontentsline{toc}{section}{Executive summary}
A TolTEC detector timestream is placed on a top-of-atmosphere,
point-source-equivalent flux-density scale by applying one detector-specific
calibration factor from the selected observation-specific array property table
(APT), whose calibration descends from Beammap calibration evidence, and one
atmospheric correction for the target sample. For detector (d), time
(t), observation (o), and array (a(d)), the organizing equation is
\begin{equation}
 d^{\rm cal}_d(t)=F^{\rm APT}_{od}\,
 C_{a(d),\alpha}\!\left[\tau_{225,o}(t),EL_o(t)\right]x^{xs}_d(t).
 \label{eq:science-central}
\end{equation}
Here \(x^{xs}=\Delta f/f_{\rm res}\) is the dimensionless ordinary \code{xs}
calibration
input, \(F^{\rm APT}\) is the selected row's \code{flxscale}, and \(C_a\) is
the inverse atmospheric transmission for the target observation. The result
is a point-source-equivalent, beam-peak-normalized amplitude intended to
produce \(\mjybeam\) when the downstream map and filtering response has unit
peak, or has been explicitly renormalized to unit peak.

The quantity \(\tau_{225}\) is dimensionless zenith optical depth measured at
225 GHz, and \(X\) is the sample airmass. Their broadband, array-dependent
mapping is discussed in Section~\ref{sec:atmosphere}.

The selected \code{flxscale} is applied exactly once. If an explicitly
approved TolProj child transformation has already embodied a pointing-derived
correction in the selected child APT, that ancestry is recorded but the
correction is not multiplied again. None of \code{responsivity}, \code{sens},
a parent \code{flxscale}, or an opaque \code{fcf} is an additional absolute
signal multiplier.

\StatusBox{\begin{tabularx}{\linewidth}{@{}L{0.28\linewidth}Y@{}}
\textbf{Intended science model} & Equation~\eqref{eq:science-central}, with one selected detector factor and the target sample's atmospheric correction.\\
\textbf{Numerical atmosphere status} & Content-bound operator \code{am12_fixed_djf25_piecewise_linear_los_tau_v1}; support \(0\le\tau_{225}\le0.25\), \(25^\circ\le EL\le80^\circ\), and \(\alpha\in\{-1,0,2,4\}\); fail closed outside.\\
\textbf{Implementation status} & Not assessed. This rationale neither audits nor certifies a Citlali implementation.\\
\textbf{Observational status} & Not validated here. Atmosphere fidelity, repeatability, and absolute source recovery require separate evidence.\\
\end{tabularx}}

Measurement-noise estimation is downstream of CAL. CAL should ultimately
report array-dependent systematic calibration uncertainty; the required
producer mechanisms and covariance model are not yet available. Common terms
do not average down as independent detector samples. Formal machine states,
identity mechanics, exact requirements, and pathological tests remain in the
engineering contract and are traced in the appendices.

\tableofcontents\clearpage

\section{What is being calibrated?}\label{sec:input}
SCI-CAL v0.1 is restricted to the ordinary \code{xs} detector stream. The
selected observation-specific APT and its associated calibration lineage
identify the source-calibration row used for each target detector. An
ambiguous or unsupported match produces no calibrated value for that detector.

Ordinary \code{xs} is the fitted KID resonant-frequency fractional change
\(\Delta f/f_{\rm res}\). It is dimensionless, and increasing \code{xs} means
increasing absorbed optical power. The fitted resonance is expected to match
the active probe tone. The stream responds to changes in total optical loading
but has no absolute DC sensitivity. SCI-CAL performs no additive baseline
subtraction or additional normalization.

\DecisionBox{Input signal and operating domain}{Ordinary \code{xs} has the
definition above and is treated as linear over expected operational loadings.
A wrong Tune invalidates that premise. Displacement beyond about half a
resonance width is a high-noise quality condition, not an alternate
calibration observable.}

Polarization calibration, auxiliary measured streams, surface-brightness or
temperature calibration, extended-source calibration, and integrated
photometry do not inherit Equation~\eqref{eq:science-central}; they are outside
v0.1.

\section{The physical calibration model and reduction chain}\label{sec:model}
A minimal generative picture separates detector response from atmospheric
attenuation. For a point-source-equivalent top-of-atmosphere signal \(s_d(t)\),
the input can be represented schematically as
\begin{equation}
 d^{\rm in}_d(t)=R_{od}T_{a(d)}(t)s_d(t)+n_d(t),
 \label{eq:generative}
\end{equation}
where \(R\) is an abstract detector response in input units per \(\mjybeam\),
\(T_a\) is transmission, and \(n\) is measurement noise. The central equation
is the inverse operation when \(F^{\rm APT}=R^{-1}\) and \(C_a=T_a^{-1}\).
Here \(R\) is not the APT field \code{responsivity}. This motivates factor
direction; the frozen SCI-BEAM contract establishes the source-factor
derivation described below.

\begin{figure}[ht]\centering
\flowboxsmall{Calibrator model\\and epoch}\(\longrightarrow\)
\flowboxsmall{Beammap calibration\\source atmosphere, amplitude, beam, pointing}\(\longrightarrow\)
\flowboxsmall{Source APT\\\code{flxscale}}\(\longrightarrow\)
\flowboxsmall{TolProj-selected\\observation-specific child APT}\(\longrightarrow\)
\flowboxsmall{TolTECA delivery\\exact immutable child APT}\\[0.08in]
\makebox[\linewidth][r]{\(\downarrow\)\hspace{0.08\linewidth}}\\[-0.02in]
\begin{tabular}{@{}c@{\ }c@{\ }c@{\ }c@{\ }c@{}}
\flowbox{Target \code{xs} +\\target atmosphere} & \(\longrightarrow\) &
\flowbox{SCI-CAL\\apply once} & \(\longleftarrow\) &
\flowbox{Selected child APT\\\code{flxscale}}\\[-0.02in]
\multicolumn{5}{c}{\(\downarrow\)}\\[-0.02in]
\multicolumn{5}{c}{\flowbox{Calibrated\\timestream}\(\longrightarrow\)
\flowbox{MAP/FLT realized\\response}}
\end{tabular}
\caption{Scientific lineage from calibrator evidence to target-observation calibration. Beammap calibration produces the source APT factor; TolProj selects or creates the observation-specific child APT; TolTECA delivers that exact immutable artifact without changing its value or meaning; SCI-CAL consumes its selected \code{flxscale} and applies the target atmosphere once.}\label{fig:lineage}
\end{figure}

SCI-CAL lies between upstream definition of the signal, APT, opacity, and
airmass and downstream PTC cleaning, filtering, weighting, mapmaking, and
inference. The
relevant ordering cases are
\begin{align}
L_d[F_d x_d]&=F_d L_d[x_d], &&\text{fixed scalar, single-detector linear operator},\nonumber\\
P[F\boldsymbol x]&\ne F P[\boldsymbol x], &&\text{detector-mixing operator in general},\nonumber\\
L_d[C_d(t)x_d(t)]&\ne C_d(t)L_d[x_d(t)], &&\text{sample-dependent atmosphere in general}.
\label{eq:noncommute}
\end{align}
SCI-CAL applies the absolute and target-atmosphere factors before PTC. PTC then
owns DC removal and correlated atmospheric/common-mode cleaning. Common-mode
removal and PCA mix detectors; a PCA basis may also be data-dependent, so the
declared order is scientifically material. Downstream uncertainty or response
companions on the calibrated domain must use the same realized multiplier and
support.

\section{From a calibrator Beammap to the target observation}\label{sec:flxscale}
Beammap analysis measures detector response to a calibrator whose
top-of-atmosphere flux model is adopted for the relevant epoch and photometric
convention. The frozen SCI-BEAM authority defines
\begin{equation}
 F^{\rm source}_{sd}=
 \frac{F_{\rm TOA,nom}}
      {A_{\rm TOA\mbox{-}equivalent,sd}}.
 \label{eq:flxscale-concept}
\end{equation}
Its units are \(\mjybeam/(\Delta f/f)\). Beammap fits the standardized detector
response with the declared finite-source-convolved effective-core model. Its
generating record identifies the calibrator and epoch, source model, nominal
beam, passband/spectrum, source-atmosphere treatment, fit support, acceptance,
and available uncertainty/covariance.

Beammap calibration and source-APT production own the calibrator model, the
source-observation atmosphere, the Beammap amplitude and beam/template fit,
the pointing treatment in source calibration, and the derivation and meaning
of source-APT \code{flxscale}. TolProj owns target-to-source detector
association, selection or creation of the observation-specific child APT, and
only explicitly approved child transformations. TolTECA owns delivery of the
exact immutable selected child artifact and its identity; delivery changes
neither the child value nor its producer- and transformer-owned scientific
meaning. SCI-CAL consumes that delivered child row as the value boundary,
applies its \code{flxscale} once, and applies target atmosphere. The downstream
MAP/FLT response owner owns the realized mapmaker, gridding, kernel, and filter
response.

This follows the producer--transformer--delivery--consumer rule: the producer
owns a quantity's meaning, a transformer owns only its approved transformation
and lineage, delivery preserves the exact selected artifact, and a consumer
applies the selected value without reinterpretation.
If
\(F^{\rm child}=P F^{\rm parent}\), the right side records ancestry; neither
\(P\) nor \(F^{\rm parent}\) is then applied to the target signal.

\begin{table}[ht]\centering\small
\caption{Roles of calibration-related quantities.}\label{tab:factor-roles}
\begin{tabularx}{\linewidth}{@{}L{0.17\linewidth}L{0.27\linewidth}L{0.18\linewidth}L{0.12\linewidth}Y@{}}
\toprule Quantity & Physical role & Source & Signal factor? & Correlation or uncertainty scope \\
\midrule
\code{flxscale} & Selected absolute detector scale & Selected observation-specific child APT; derived from source-APT/Beammap calibration & Yes, once & Detector plus shared calibrator ancestry \\
Pointing correction & Correction in APT ancestry & Beammap source calibration or an explicitly approved child transformation & No, if embodied & Detector/observation correlated \\
\code{responsivity} & Relative response or diagnostic & Selected APT lineage & No & Transfer role unresolved \\
\code{sens} & Sensitivity or approximate weight & Beammap analysis & No & Detector/cohort; not total error \\
Target \(C_a\) & Inverse transmission for target sample & Opacity, airmass, array operator & Yes, once & Time, observation, and array correlated \\
Compatibility \code{fcf} & Declared decomposition of known roles & Compatibility boundary & Not independently & Inherits constituent scopes \\
\bottomrule\end{tabularx}\end{table}

Every target detector needs one unique, traceable, observation-valid match to
one measured APT row. Exact mechanics remain in the engineering contract; an
unmatched or ambiguous detector is excluded rather than calibrated from row
position or a plausible default.

\DecisionBox{Calibration selection and optional rescale}{The default is the
temporally closest scientifically accepted Beammap source APT; no universal
maximum age or same-Tune, loading, focus, or detector-state condition is
imposed. At the reducing scientist's direction, TolProj may create a child APT
with one per-array photometric rescale based on a suitable recent ALMA, SMA,
planet, or well-calibrated asteroid reference. Its lineage records the source,
authorities/epochs, inferred and reference fluxes, spectral convention,
factor, and uncertainty when available. Primary-mirror efficiency variation
with aberration remains a hypothesis, not an established calibration law.}

\section{Atmospheric extinction correction}\label{sec:atmosphere}
The LMT water-vapor radiometer supplies \(\tau_{225}\). Its raw readings are
written into the observation's \code{tel*.nc} files approximately every five
minutes and are interpolated in time to the detector sample grid. Full sample
airmass is used with \(X_{\rm ref}=0\). The retained operator is
\code{am12_fixed_djf25_piecewise_linear_los_tau_v1}, generated from the fixed
\code{LMT_DJF_25.amc} profile.

For the frozen TolTECA-v1 array response \(R_a(\nu)\), reference spectrum
\(S_\alpha\propto\nu^\alpha\), and AM transmission, it evaluates
\begin{equation}
 T^{\rm eff}_{a,\alpha}(\tau,EL)=
 \frac{\int R_a S_\alpha T_{\rm AM}\,d\nu}
      {\int R_a S_\alpha\,d\nu},\qquad
 C_{a,\alpha}=1/T^{\rm eff}_{a,\alpha}.
 \label{eq:atmos-essential}
\end{equation}
The supplied energy-weighted throughput is integrated by composite trapezoid
on its exact nodes without spectral extrapolation. The authoritative ordinate
is \(-\ln T^{\rm eff}\): PCHIP in elevation at each opacity anchor, then
piecewise-linear interpolation across opacity anchors. The surface has an
analytic unity anchor at zero and no switch at \(\tau_{225}=0.15\).

\DecisionBox{Authoritative atmosphere operator}{The exact support is
\(0\le\tau_{225}\le0.25\), \(25^\circ\le EL\le80^\circ\), and
\(\alpha\in\{-1,0,2,4\}\), with default \(\alpha=0\) and no interpolation or
extrapolation in \(\alpha\). The operator fails closed outside support. Its
authority commit is \code{7156881bd1a47e8cece97b8c541a013c93ac03e1}; the
contract and node-table digests are recorded in the engineering authority.}

CAL assigns one class to the entire observation and never splits it into
science and engineering pieces. The nominal \(0.15\) science and approximately
\(0.25\) engineering limits are flexible experience-based operational
guidance, not physical or calibration-error discontinuities. A 0.025 tolerance
allows momentary excursions without invalidating the whole observation. This
quality tolerance never expands the numerical operator: an individual sample
outside support remains unavailable. Policy changes require owner approval.

\section{Photometric convention and downstream response}\label{sec:response}
The intended output is a top-of-atmosphere, point-source-equivalent,
beam-peak-normalized amplitude. ``Beam'' denotes the originating Beammap
response normalization; it is neither pixel area nor an approved definition of
surface brightness or integrated flux.

For a source whose flux density is \(S\) mJy, originating unit-peak template
\(P_B\), and realized mapmaker/gridding/filter response (L),
\begin{equation}
 m_{\rm peak}=\rho_L S\ \mjybeam,\qquad
 \rho_L=[L P_B](\boldsymbol 0).
 \label{eq:response-main}
\end{equation}
The literal peak equals \(S\) only for \(\rho_L=1\), or after documented
renormalization. A valid smoothing kernel can broaden a point source and
reduce its peak; retaining the old label without response correction would
then misstate flux.

TolTEC is broadband. The per-array reference frequencies used by
\code{toltec_beammap} are 272, 214, and 150 GHz for a1100, a1400, and a2000.
Calibration-source spectra are source dependent and commonly represented by
\(S_\nu\propto\nu^\alpha\). BEAM/\code{toltec_beammap} supplies one per-array
\(\mjybeam\) input to \code{flxscale}; SCI-CAL applies it without a second
source color correction.

\DecisionBox{Broadband photometric convention}{Target-source color correction
is a downstream scientist action. Until that action, the product retains its
declared reference-spectrum convention. The atmosphere model uses one
array-average passband per array; detector/network passband variation and its
uncertainty are presently unavailable.}

\section{Uncertainty: measurement noise is not total calibration error}\label{sec:uncertainty}
Measurement-noise estimation is downstream of CAL; CAL does not create a
statistical variance or weight. For fixed calibration, atmosphere, detector
association, and response, let \(M_j=F^{\rm APT}_{od}C_{a(d),\alpha}(t)\) and
\(A=\operatorname{diag}(M_j)\). A downstream covariance on the same support must
remain consistent with
\begin{equation}
 C_{\rm cal\mid\theta}=A C_{\rm in\mid\theta}A^{\mathsf T},\qquad
 \sigma^2_{{\rm cal},j}=M_j^2\sigma^2_{{\rm in},j},\quad
 w_{{\rm cal},j}=w_{{\rm in},j}/M_j^2.
 \label{eq:conditional-main}
\end{equation}
This is a downstream consistency identity, not a CAL uncertainty product.
Missing uncertainty is unavailable, never zero. Calibration-state uncertainty
is different. A fractional scale
uncertainty (s) common to affected samples contributes
\begin{equation}C_{\rm common}=s^2\boldsymbol\mu\boldsymbol\mu^{\mathsf T},\label{eq:common-main}\end{equation}
which does not decrease merely because more identically affected samples are
averaged.

\begin{table}[ht]\centering\small
\caption{Science-facing uncertainty budget. ``Not established'' means the frozen authority does not state a present numerical product.}\label{tab:uncertainty-budget}
\begin{tabularx}{\linewidth}{@{}L{0.21\linewidth}L{0.24\linewidth}L{0.13\linewidth}L{0.15\linewidth}Y@{}}
\toprule Quantity & Typical correlation scope & In sample weight? & Current status & Interpretation \\
\midrule
Measurement noise & Sample/detector, possibly correlated & Downstream & Downstream-owned & Must obey \(A C A^T\) on calibrated support \\
Detector \code{flxscale} fit & Detector/Beammap cohort & No & Unavailable; BEAM mechanism needed & Calibration systematic \\
Calibrator scale/model & Observations and/or band & No & Not established & Does not average down \\
Pointing rescale & Array/observation & No & Child-APT mechanism undefined & Embodied once \\
Atmosphere, WVR, operator & Observation-common driver; array response & No & Uncertainty unavailable & CAL must ultimately estimate \\
Passband/source spectrum & Array & No & Variation unavailable & Broadband systematic \\
Beam/filter response & Product/reduction state & No & Downstream-owned & Needed for literal peak \\
\code{sens} approximation & Detector/cohort & Only if declared & Not established & Never total error \\
\bottomrule\end{tabularx}\end{table}

CAL should ultimately report an array-dependent systematic calibration error
for scientist use. Scale terms are common within each array. WVR state is
observation-common across arrays but has an array-dependent response;
numerical cross-array covariance is unavailable. Until the source-APT,
TolProj-rescale, WVR, and covariance mechanisms exist, total uncertainty and
significance remain unavailable, not zero.

\section{Validity limits and what a science user sees}\label{sec:validity}
\begin{longtable}{L{0.24\linewidth}L{0.14\linewidth}L{0.14\linewidth}L{0.39\linewidth}}
\caption{Meaning of calibration status for science users.}\label{tab:user-states}\\
\toprule Status & Calibrated value? & Conditional weight? & Interpretation and action \\
\midrule\endfirsthead
\toprule Status & Calibrated value? & Conditional weight? & Interpretation and action \\
\midrule\endhead
Structurally complete within nominal domain & Yes, once numeric inputs are supported by approved definitions & If supplied & Algebra and inputs complete; atmosphere fidelity and on-sky performance remain separate. \\
Calibration disabled & No & No calibrated weight & Intentional uncalibrated mode; do not relabel input. \\
Unsupported stream or unit & No & No & A separate scientific contract is required. \\
Invalid/missing atmosphere & No & No & Correct the owning input or exclude segment. \\
Outside operator support & No & No & No extrapolation, clamp, or legacy fallback. \\
Engineering-only observation & Yes on operator-supported samples & Downstream & Same calibration operator; no science-quality claim. \\
Conditional uncertainty unavailable & Signal may exist & No & Do not interpret missing uncertainty as zero. \\
Total uncertainty incomplete & Signal may exist & Conditional only & Report noise separately from systematics. \\
Response basis incomplete & Signal may exist & Conditional only & Withhold literal peak meaning. \\
\bottomrule\end{longtable}
The engineering state \code{science-qualification-eligible} is rendered here
as ``structurally complete but not yet scientifically validated.''

\section{What has to be validated, and what each validation establishes}\label{sec:validation}
\begin{enumerate}[label=\textbf{\arabic*.},leftmargin=*]
\item \textbf{Structural calibration correctness.} Factor direction and
once-only challenges; unit and support checks; unique detector to APT association;
zero opacity, node, seam, monotonicity, and no extrapolation behavior;
downstream covariance consistency; and response bookkeeping. Software-oriented
permutation, duplicate-key, replay, sentinel, and application-count tests
remain in the engineering contract.
\item \textbf{Atmosphere-model representation fidelity.} Compare the finite
operator with its named higher-fidelity generating calculation over each
array's declared opacity, elevation, spectral, and passband domain. Report the
actual result; the prior approximately one-percent figure is only a benchmark.
\item \textbf{Beammap closure.} Reduce calibrated Beammaps to source APTs,
validate them with \code{toltec_beammap}, enter accepted APTs into
\code{apt_library}, associate them through TolProj, and reduce each generating
Beammap through the ordinary science path. Recovery of the input calibrator
flux establishes APT--TolProj--CAL--downstream closure, not independent
absolute calibration accuracy.
\item \textbf{Associated-pointing transfer.} Apply each accepted Beammap APT
to its associated pointing observation through the same TolProj and science
path. When the pointing source has an adequate independent ALMA, SMA, planet,
or well-calibrated asteroid flux, this is a stronger calibration-transfer
test. Repeat over the observations actually available and report separately
for every array.
\end{enumerate}
The record identifies observations, source and child APTs, TolProj
associations, flux references, arrays, atmosphere, recovered/input ratios or
residuals, repeatability, absolute recovery, and observed condition
dependence. No arbitrary combinatorial matrix or minimum sample size is
required. The one-percent, five-percent, and five-to-ten-percent figures are
reporting benchmarks, not automatic pass/fail ceilings. The task is to report
what the system achieves honestly and as accurately as possible; final
scientific acceptance is an explicit owner decision.

\section{Scientific owner decisions and remaining unavailable products}\label{sec:decisions}
\begin{longtable}{L{0.06\linewidth}L{0.22\linewidth}L{0.32\linewidth}L{0.31\linewidth}}
\caption{Decision register for science-rationale revision r0.5.}\label{tab:decisions}\\
\toprule ID & Issue & Owner disposition & Remaining limitation \\
\midrule\endfirsthead
\toprule ID & Issue & Owner disposition & Remaining limitation \\
\midrule\endhead
Q01 & Physical \code{xs} definition & Decided: dimensionless \(\Delta f/f_{\rm res}\), positive optical-power direction, no CAL baseline & Wrong Tune invalidates; half-width excursion is high-noise state \\
Q02 & Baseline and ordering & Decided: CAL before PTC; PTC owns DC/common-mode removal & Downstream transfer remains separately measured \\
Q03 & \code{flxscale} derivation & Decided by frozen SCI-BEAM authority & Producer uncertainty presently unavailable \\
Q04 & Calibrator-target transfer & Decided: closest accepted Beammap; optional scientist-directed per-array TolProj rescale & Rescale uncertainty mechanism unavailable \\
Q05 & Broadband convention & Decided: 272/214/150 GHz, source-dependent spectrum, downstream target color correction & Detector/network passband variation unavailable \\
Q06 & Atmosphere operator & Decided: content-bound WVR/AM/passband operator & WVR/model uncertainty unavailable \\
Q07 & Opacity/observation policy & Decided: one CAL class, 0.025 excursion tolerance, owner changes policy & Operational guidance is not a performance claim \\
Q08 & Numerical uncertainty & Decided ownership and correlation scope & Total uncertainty/significance unavailable \\
Q09 & Criteria for validated science use & Decided closure-and-transfer workflow and reporting fields & Final acceptance remains an owner decision \\
\bottomrule\end{longtable}

\section*{References}\addcontentsline{toc}{section}{References}
\begin{enumerate}[label={[\arabic*]},leftmargin=*]
\item \emph{SCI-CAL Detector Calibration, Atmospheric Extinction, and Signal Transfer Scope Brief}, owner-approved 2026-08-16.
\item \emph{SCI-CAL Independent Scientific Core}, with v0.1 supersession cover.
\item \emph{SCI-CAL Sanitized Conventions and Ownership}, owner-approved 2026-08-16.
\item \emph{TolTEC v1 Modeled Array Passband Authority Manifest}, bounded reference with unresolved semantics.
\item \emph{SCI-CAL Scientific Owner Decisions for r0.5/r0.4}, approved 2026-08-20.
\item \emph{SCI-CAL Atmosphere Operator Machine Contract}, authority commit\\
\code{7156881bd1a47e8cece97b8c541a013c93ac03e1}.
\end{enumerate}
The sanitized conventions/ownership reference records the lineage of its
approved extract. No broader source in that lineage was independently opened
or used as an author input.
The r0.5 owner-decision record supplies the approved \code{xs}, pipeline,
transfer, photometric, quality, uncertainty, and validation dispositions. The
content-bound machine contract supplies the exact numerical atmosphere
authority; the frozen SCI-BEAM contract remains the source-factor authority.

\clearpage\appendix
\section{Notation and formal derivations}\label{app:formal}
\begin{longtable}{L{0.20\linewidth}L{0.70\linewidth}}
\caption{Notation retained for formal traceability.}\\
\toprule Symbol & Meaning \\
\midrule\endfirsthead\toprule Symbol & Meaning \\
\midrule\endhead
(o,d,t,a(d)) & Observation, detector, sample time, detector array \\
\(d^{\rm in},d^{\rm cal}\) & Upstream-defined input and calibrated output \\
\(F^{\rm APT}_{od}\) & Selected row's oriented absolute \code{flxscale} \\
\(\tau_{225},X,\ell\) & Zenith opacity, airmass, line-of-sight coordinate \\
\(\lambda_{a,\alpha},T_{a,\alpha},C_{a,\alpha}\) & Interpolated log-transmission ordinate, transmission, correction \\
\(M_j,A\) & Realized scalar multiplier and diagonal operator \\
\(C_{\rm in},C_{\rm cal},C_\theta\) & Conditional input/output and nuisance covariance \\
\(P_B,L,\rho_L\) & Beammap template, downstream response, peak response \\
\bottomrule\end{longtable}
The owner-declared CAL input has no additive operation:
\begin{equation}d^{\rm in}_j=x^{xs}_j.\label{eq:appendix-affine}\end{equation}
At each opacity anchor the authoritative \(\lambda=-\ln T^{\rm eff}\) surface
uses PCHIP in elevation. Between opacity anchors,
\begin{align}
u&=\frac{\tau-\tau_k}{\tau_{k+1}-\tau_k},\\
\lambda_{a,\alpha}(\tau,EL)&=(1-u)\lambda_{a,\alpha}(\tau_k,EL)
+u\lambda_{a,\alpha}(\tau_{k+1},EL),\\
(T_{a,\alpha},C_{a,\alpha})&=(e^{-\lambda_{a,\alpha}},e^{\lambda_{a,\alpha}}).
\label{eq:appendix-atmos}
\end{align}
The declared log-transmission ordinate is interpolated before exponentiation. Exact nodes are
recovered, seams continuous, \(T_a(0)=C_a(0)=1\), transmission decreases,
correction increases, and extrapolation is unauthorized.

For fixed state \(\theta\),
\begin{equation}
\mathbb E[\boldsymbol d^{\rm cal}\mid\theta]=A\mathbb E[\boldsymbol d^{\rm in}\mid\theta],\qquad C_{\rm cal\mid\theta}=A C_{\rm in\mid\theta}A^{\mathsf T}.
\end{equation}
For small nuisance uncertainty with
\(H_{jr}=\partial\ln|M_j|/\partial\theta_r\),
\begin{equation}C_{\rm nuis}\simeq\operatorname{diag}(\boldsymbol\mu)H C_\theta H^{\mathsf T}\operatorname{diag}(\boldsymbol\mu).\end{equation}
Shared-data calibration requires material cross covariances. Large,
asymmetric, seam-spanning, or discrete uncertainty
requires nonlinear samples, a joint posterior, or bounded sensitivity study.

\section{Exact machine states and formal conformance routing}\label{app:states}
The engineering authority distinguishes requested, effective,
observation-resolved, and realized states and valid signal. Its exact named
states include:
\begin{itemize}
\item \code{disabled/uncalibrated};
\item \code{outside-supported-calibration};
\item \code{engineering-only}; and
\item \code{science-qualification-eligible}.
\end{itemize}
It separately records conditional-uncertainty availability,
total-uncertainty completeness, and response completeness.
Table~\ref{tab:user-states} translates these for users; it does not replace
them. Appendix~\ref{app:crosswalk} maps every formal item to this rationale and
its continuing engineering home.

\section{Provenance and document authority}\label{app:provenance}
The approved input content digests are:
\begin{itemize}
\item Scope Brief: {\footnotesize\code{a3d6332dcc98bfbb638a45aea9be830edca693a4fe7e65d831b5880603c16578}};
\item independent core: {\footnotesize\code{106755520b048f601bc60fd04e7b6020e6fa470480ac3105fa7ba269c730a4fe}};
\item passband manifest: {\footnotesize\code{2756908181cc466550399ec0a869e6671de7912bd3a935f9aeebf63e3e826617}};
\item conventions/ownership: {\footnotesize\code{7e9a630fd183ca04bc3d8bbd21b5e801776b9aad7bd084d48fa7f2c572766520}}.
\end{itemize}
The engineering contract remains normative for requirements 001--050,
assumptions 001--011, and edge predictions 001--030. Rationale revision r0.5
changes genre and presentation, not the governed v0.1 algebra.

\section{Crosswalk to the engineering contract}\label{app:crosswalk}
\small
The engineering contract remains normative. This compact map locates every
formal assumption, requirement, and edge prediction without reproducing its
full audit-oriented text in the science narrative.

\begin{longtable}{L{0.30\linewidth}L{0.26\linewidth}L{0.35\linewidth}}
\caption{Assumption routing.}\\
\toprule Formal IDs & Scientist-facing location & Continuing engineering home \\
\midrule\endfirsthead\toprule Formal IDs & Scientist-facing location & Continuing engineering home \\
\midrule\endhead
\Assump{001}, \Assump{002} & Sections~\ref{sec:input}--\ref{sec:model}; Q01--Q02 & Approved input support, ordinary channel, and binding \\
\Assump{003} & Section~\ref{sec:flxscale}; Figure~\ref{fig:lineage}; Q03--Q04 & Source APT, TolProj-selected child APT, immutable TolTECA delivery, and association \\
\Assump{004} & Section~\ref{sec:input}; Appendix~\ref{app:formal} & No-additive-operation CAL input convention \\
\Assump{005} & Sections~\ref{sec:model}--\ref{sec:flxscale}; Q03--Q05 & Complete factor lineage, direction, plane, and once-only use \\
\Assump{006} & Section~\ref{sec:atmosphere}; Q06--Q07 & WVR inputs, time support, operator, and observation policy \\
\Assump{007} & Section~\ref{sec:uncertainty} & Downstream conditional-state covariance boundary \\
\Assump{008} & Section~\ref{sec:response} & Originating and downstream response ownership \\
\Assump{009} & Sections~\ref{sec:atmosphere}, \ref{sec:validity}; Q07 & Whole-observation quality class \\
\Assump{010} & Sections~\ref{sec:atmosphere}--\ref{sec:response}; Q05 & Passband limitations and photometric compatibility \\
\Assump{011} & Sections~\ref{sec:atmosphere}, \ref{sec:decisions}; Q06 & Content-bound numeric operator authority \\
\bottomrule
\end{longtable}

\begin{longtable}{L{0.30\linewidth}L{0.26\linewidth}L{0.35\linewidth}}
\caption{Requirement routing.}\\
\toprule Formal IDs & Scientist-facing location & Continuing engineering home \\
\midrule\endfirsthead\toprule Formal IDs & Scientist-facing location & Continuing engineering home \\
\midrule\endhead
\Req{001}, \Req{002}, \Req{003}, \Req{004}, \Req{005} & Sections~\ref{sec:input}, \ref{sec:validity}; Q01--Q09 & Scope, typed state, authority snapshot, validity, and one-way resolution \\
\Req{006}, \Req{007}, \Req{008}, \Req{009}, \Req{010}, \Req{011}, \Req{012} & Section~\ref{sec:flxscale}; Figure~\ref{fig:lineage} & Acquisition/source/child identity, binding, association, delivery, and lifetime \\
\Req{013}, \Req{014}, \Req{015}, \Req{016} & Sections~\ref{sec:model}--\ref{sec:flxscale}; Q03--Q04 & SCI-BEAM generating record, optional child rescale, and once-only multiplier \\
\Req{017}, \Req{018}, \Req{019}, \Req{020} & Section~\ref{sec:flxscale}; Table~\ref{tab:factor-roles} & Relative roles, compatibility, producer/transformer/delivery/consumer lineage \\
\Req{021}, \Req{022}, \Req{023}, \Req{024}, \Req{025}, \Req{026} & Section~\ref{sec:atmosphere}; Appendix~\ref{app:formal} & WVR meaning, temporal/operator support, and invariants \\
\Req{027}, \Req{028}, \Req{029}, \Req{030}, \Req{031} & Sections~\ref{sec:atmosphere}, \ref{sec:validity}; Q05--Q07 & Whole-observation policy, sample support, passband, and spectrum \\
\Req{032}, \Req{033}, \Req{034} & Section~\ref{sec:uncertainty} & Downstream covariance/weight boundary and unavailable state \\
\Req{035}, \Req{036}, \Req{037}, \Req{038}, \Req{039} & Sections~\ref{sec:model}, \ref{sec:uncertainty}; Appendix~\ref{app:formal}; Q08 & Systematic scopes, unavailable covariance, and companions \\
\Req{040}, \Req{041}, \Req{042}, \Req{043}, \Req{044} & Sections~\ref{sec:flxscale}, \ref{sec:response} & Producer-owned beam/template basis, realized response, and unit exclusions \\
\Req{045}, \Req{046}, \Req{047}, \Req{048} & Sections~\ref{sec:flxscale}, \ref{sec:validity}; Appendix~\ref{app:states}; Q01--Q09 & Claim-specific states, reasons, complete lineage, and product links \\
\Req{049}, \Req{050} & Sections~\ref{sec:response}, \ref{sec:validation}; Q09 & Separate closure, transfer, fidelity, performance, and acceptance claims \\
\bottomrule
\end{longtable}

\begin{longtable}{L{0.30\linewidth}L{0.26\linewidth}L{0.35\linewidth}}
\caption{Edge-prediction routing.}\\
\toprule Formal IDs & Scientist-facing location & Continuing engineering home \\
\midrule\endfirsthead\toprule Formal IDs & Scientist-facing location & Continuing engineering home \\
\midrule\endhead
\Edge{001}, \Edge{002}, \Edge{003}, \Edge{004}, \Edge{005} & Section~\ref{sec:atmosphere}; Appendix~\ref{app:formal} & Zero, monotonicity, zenith, exact-node, and seam checks \\
\Edge{006}, \Edge{007}, \Edge{008}, \Edge{009}, \Edge{010} & Sections~\ref{sec:atmosphere}, \ref{sec:validity} & Domain, temporal bracketing, engineering and outside-support states \\
\Edge{011}, \Edge{012}, \Edge{013} & Section~\ref{sec:flxscale} & Unity, omitted/duplicate/inverted, corrected-APT challenges \\
\Edge{014}, \Edge{015}, \Edge{016}, \Edge{017}, \Edge{018}, \Edge{019} & Section~\ref{sec:flxscale} & Row/order/network, missing/duplicate, association/design tests \\
\Edge{020}, \Edge{021}, \Edge{022}, \Edge{023}, \Edge{024} & Section~\ref{sec:uncertainty} & Downstream consistency, unavailable, common, and nonlinear uncertainty tests \\
\Edge{025}, \Edge{026}, \Edge{027} & Section~\ref{sec:response} & Unit-peak source, response-changing kernel, beam approximation \\
\Edge{028} & Sections~\ref{sec:input}, \ref{sec:validity} & Unsupported stream or target \\
\Edge{029} & Section~\ref{sec:input}; Appendix~\ref{app:formal} & Zero signal and no-additive-operation convention \\
\Edge{030} & Sections~\ref{sec:flxscale}, \ref{sec:validity} & Observation-order state replay \\
\bottomrule
\end{longtable}

\end{document}
